\documentclass[10pt]{article}

% Packages

\usepackage[
    a4paper,
    top=2.5cm,
    bottom=2.5cm,
    left=3cm,
    right=2.5cm,
    marginparwidth=1.75cm
]{geometry}

\usepackage{subfiles}

% Spacing packages
\usepackage{setspace}
\usepackage{parskip}

% Extra formatting packages
% in Main.tex, Präambel:
\usepackage{csquotes}
\newcommand{\game}[1]{\textit{#1}}   % Spieltitel
\newcommand{\code}[1]{\texttt{#1}}   % Klassen-,
\usepackage{enumitem}

% Language packages
\usepackage[T1]{fontenc}
\usepackage[ngerman]{babel}

% Formalities
\usepackage[backend=biber, style=authortitle-ibid]{biblatex}

% Graphical packages
\usepackage{graphicx}
\usepackage{float}
\usepackage{abstract}

% Mathematical packages
\usepackage{amssymb}
\usepackage{amsmath}
\usepackage{mathtools}
\usepackage{siunitx}
\usepackage{caption}
\usepackage{array}
\usepackage{longtable}

% Definitions and setup

\onehalfspacing
\addbibresource{Source/Sources.bib}

\setlength{\absleftindent}{0pt}
\setlength{\absrightindent}{0pt}

\graphicspath{
    { ./Assets/ }
}

% Document

\begin{document}

% Title page

\begin{titlepage}
    \begin{center}
        \vspace*{2cm}

        \textbf{\Huge{Project Polearm}} \\ [2em]
        \textbf{\Large Die Entwicklung eines Hack-and-Slash-Videospiels} \\ [0.6em]
        \large{Konzeption, Umsetzung und Evaluation von Gegnerverhalten und Leveldesign in einem Hack-and-Slash-Prototyp} \\ [2em]

        \Large{Rodin Arslan, Zhile Zhu} \\ [0.5em]
        \large{Klasse G23g}

        \vspace{3cm}

        \LARGE{\textbf{Maturitätsarbeit}} \\ [0.3em]
        \Large{Begleitbericht} \\ [0.5em]
        \large{Kantonsschule Baden}

        \vfill

        \Large{\textbf{Betreuung und Bewertung}} \\ [0.5em]
        \large{Betreuer: Michael Schneider, Kantonsschule Baden, Dept. Informatik} \\
        \large{Zweitbewerter: Dr. Christophe Bersier, Kantonsschule Baden, Dept. Physik} \\

        \vspace{2cm}
        \large{10. November 2026}
    \end{center}
\end{titlepage}

\newpage

% Abstract

\begin{abstract}
    Die vorliegende Maturitätsarbeit beschäftigt sich mit der Erkenntnisfrage, wie Gegnerverhalten und Gegnerplatzierung in einem 3D-Hack-and-Slash-Prototyp konzipiert sein müssen, damit ein Kampfabschnitt mit nachvollziehbarem Spielfluss, angemessener Schwierigkeit und befriedigendem Spielgefühl entsteht. Des Zweck besteht darin, selbst einen solchen Prototyp zu erstellen und unsere erkenntnisse zu validieren.

    Zur Konzipierung des eigenen Prototyps wurde das Gegnerverhalten in \game{Ninja Gaiden Sigma Plus} und \game{NieR: Automata} mithilfe des MDA-Frameworks analysiert. Für die Gegnerplatzierung und das Leveldesign wurden Fachliteratur und Expertenseminare ausgewertet. Aus beidem wurden acht konkrete Designvorschläge gewonnen. Diese wurden in einem Prototyp umgesetzt, der mit der Game Engine Godot entwickelt wurde. Der Prototyp hat insgesamt 12 Encounter in 11 verschiedenen Räumen und beinhalet 5 verschiedene Gegnervarianten. In Playtests mit fünf Testpersonen wurde anschliessend untersucht, wie sich diese Umsetztung auf das Spielverhalten und das Spielerlebnis auswirkte.

    \begin{itemize}[nosep, leftmargin=*]
        \item \textcolor{red}{Anzahl nennen: \enquote{acht konkrete Designvorschläge}.}
        \item \textcolor{red}{Umfang des Prototyps (ca.\ 10 Wörter): zwölf Encounter in elf Räumen, fünf Gegnertypen.}
        \item \textcolor{red}{Methode der Playtests (ca.\ 30 Wörter): Beobachtung mit lautem Denken, Pre- und Post-Questionnaire, semistrukturiertes Interview, qualitative Auswertung.}
    \end{itemize}

    Diese Playtests deuten darauf hin, dass vor allem die Begrenzung der gleichzeitig aktiven Gegner, die Kombination von Nah- und Fernkämpfern und die gestaffelte Einführung neuer Gegnervarianten zu einer angemessenen Schwierigkeit beitragen. Für den Spielfluss und das Spielgefühl konnten nur weniger Aussagekräftige tendenzen festgestellt wreden, da Faktoren ausserhalb der Gegner die aussagekraft der Ergebnisse beeinträchtigten. Vor allem von Faktoren, wie Spielercombos und unzureichender Informationsvermittlung an den Spieler ist die Rede. Aus den gewonnenen Ergebnissen konnten Verbesserungen für eine weitere Iteration der Designvorschläge und des Playtests hergeleitet werden.

    \begin{itemize}[nosep, leftmargin=*]
        \item \textcolor{red}{Konkrete Faktoren ausserhalb der Gegner nennen: träge Combos, fehlendes Signal am Ende eines Encounters (und klarstellen, dass sie Spielfluss und Spielgefühl selbst beeinflusst haben, nicht die Beobachtung).}
        \item \textcolor{red}{Vorläufige Antwort auf die Erkenntnisfrage aus Abschnitt~\ref{sec:Beantwortung} in einem Satz (ca.\ 25 Wörter).}
        \item \textcolor{red}{Einschränkung als Halbsatz: Wegen fünf Testpersonen und qualitativer Auswertung sind die Ergebnisse nur als Hinweise zu verstehen.}
    \end{itemize}
\end{abstract}

\newpage

% Table of contents

\tableofcontents

\newpage

% Main document

\subfile{Source/1Einleitung.tex}
\newpage

\subfile{Source/2TheoretischerHintergrund.tex}
\newpage

\subfile{Source/3KampfAnalyse}
\newpage

\subfile{Source/4Leveldesign.tex}
\newpage

\subfile{Source/5Umsetzung.tex}
\newpage

\subfile{Source/7Evaluation.tex}
\newpage

\subfile{Source/8Schlussfolgerung.tex}
\newpage

\subfile{Source/9Reflektion.tex}
\newpage

\subfile{Source/10Verzeichnis.tex}
\newpage

\end{document}
